% GENERATED FILE. Edit canonical lessons, then run npm run generate:books.
% canonical-source-hash: fnv1a64:481b41161297fe71
% canonical-lessons: GU-C320-yuddh, GU-C320-sena, GU-C320-sainik, GU-C320-dusman, GU-C320-akraman

\chapter{War Comes}
\label{ch:gu-war-comes}
\hlchaptermodality{320}

\begin{chapteropening}
\textbf{By the end of this chapter:} \emph{I can say who fought whom and how the war began.}

Complete the last lesson of chapter 320: I can say who fought whom and how the war began.
\end{chapteropening}

\section[\texorpdfstring{\gu{યુદ્ધ} (yuddh)}{યુદ્ધ (yuddh)}]{\gu{યુદ્ધ} (yuddh) — a war}

\label{lesson:GU-C320-yuddh}



\begin{quote}
Before the new one: say the Gujarati for a statue, then the Gujarati for to settle.
\end{quote}

\subsection*{You'll want to know: \gu{યુદ્ધ}}
\textbf{\gu{યુદ્ધ}} — \emph{yuddh} — \textquotedblleft{}a war\textquotedblright{}.

\begin{quote}
\gu{બે} \gu{રાજાઓ} \gu{વચ્ચે} \gu{યુદ્ધ} \gu{થયું}.
\end{quote}

\begin{quote}
\emph{be rājāo vacce yuddh thayũ.}
\end{quote}

\begin{quote}
War broke out between two kings.
\end{quote}

A neuter noun.

\begin{cousinweb}
\textbf{\gu{યુદ્ધ}} (\emph{yuddh}) is a tatsama, Sanskrit \emph{yuddha}, \textquotedblleft{}fought\textquotedblright{}, and as a noun \textquotedblleft{}a battle\textquotedblright{} (Monier-Williams). It is neuter: \textbf{\gu{યુદ્ધ} \gu{થયું}}, there was a war.
\end{cousinweb}

\begin{grammarlens}[title={What you've built}]
One more word for telling what happened.
\end{grammarlens}

\subsection*{Your turn}
\begin{itemize}
\raggedright
  \item \emph{Say it:} \emph{yuddh}
  \item \emph{Say it:} \emph{yuddh}, once more
  \item \emph{Say it:} \emph{be rājāo vacce yuddh thayũ.} — War broke out between two kings.
  \item \emph{Recall:} say the Gujarati for a statue, then the Gujarati for to settle, then say \emph{yuddh} again
\end{itemize}

\begin{tcolorbox}[breakable,colback=teal!4,colframe=teal!35!black,title={Before you move on}]
What does \gu{યુદ્ધ} mean? (\textquotedblleft{}A war\textquotedblright{}.) Say it once more.
\end{tcolorbox}
\section[\texorpdfstring{\gu{સેના} (senā)}{સેના (senā)}]{\gu{સેના} (senā) — an army}

\label{lesson:GU-C320-sena}



\begin{quote}
Before the new one: say the Gujarati for to settle, then the Gujarati for a war.
\end{quote}

\subsection*{You'll want to know: \gu{સેના}}
\textbf{\gu{સેના}} — \emph{senā} — \textquotedblleft{}an army\textquotedblright{}.

\begin{quote}
\gu{રાજાની} \gu{સેના} \gu{ખૂબ} \gu{મોટી} \gu{હતી}.
\end{quote}

\begin{quote}
\emph{rājānī senā khūb moṭī hatī.}
\end{quote}

\begin{quote}
The king's army was very large.
\end{quote}

A feminine noun.

\begin{cousinweb}
\textbf{\gu{સેના}} (\emph{senā}) is a tatsama, Sanskrit \emph{s\'{\={e}}nā}, \textquotedblleft{}an army\textquotedblright{} (CDIAL 13587). It is feminine: \textbf{\gu{મોટી} \gu{સેના}}, a large army.
\end{cousinweb}

\begin{grammarlens}[title={What you've built}]
One more word for telling what happened.
\end{grammarlens}

\subsection*{Your turn}
\begin{itemize}
\raggedright
  \item \emph{Say it:} \emph{senā}
  \item \emph{Say it:} \emph{senā}, once more
  \item \emph{Say it:} \emph{rājānī senā khūb moṭī hatī.} — The king's army was very large.
  \item \emph{Recall:} say the Gujarati for to settle, then the Gujarati for a war, then say \emph{senā} again
\end{itemize}

\begin{tcolorbox}[breakable,colback=teal!4,colframe=teal!35!black,title={Before you move on}]
What does \gu{સેના} mean? (\textquotedblleft{}An army\textquotedblright{}.) Say it once more.
\end{tcolorbox}
\section[\texorpdfstring{\gu{સૈનિક} (sainik)}{સૈનિક (sainik)}]{\gu{સૈનિક} (sainik) — a soldier}

\label{lesson:GU-C320-sainik}



\begin{quote}
Before the new one: say the Gujarati for a war, then the Gujarati for an army.
\end{quote}

\subsection*{You'll want to know: \gu{સૈનિક}}
\textbf{\gu{સૈનિક}} — \emph{sainik} — \textquotedblleft{}a soldier\textquotedblright{}.

\begin{quote}
\gu{સેનાનો} \gu{સૌથી} \gu{બહાદુર} \gu{સૈનિક} \gu{ઘાયલ} \gu{થયો}.
\end{quote}

\begin{quote}
\emph{senāno sauthī bahādur sainik ghāyal thayo.}
\end{quote}

\begin{quote}
The bravest soldier in the army was wounded.
\end{quote}

A masculine noun.

\begin{cousinweb}
\textbf{\gu{સૈનિક}} (\emph{sainik}) is a tatsama, Sanskrit \emph{sainika}, \textquotedblleft{}belonging to an army\textquotedblright{} (Monier-Williams), built on \emph{sēnā}, the word behind \textbf{\gu{સેના}}: one who serves in it.
\end{cousinweb}

\begin{grammarlens}[title={What you've built}]
One more word for telling what happened.
\end{grammarlens}

\subsection*{Your turn}
\begin{itemize}
\raggedright
  \item \emph{Say it:} \emph{sainik}
  \item \emph{Say it:} \emph{sainik}, once more
  \item \emph{Say it:} \emph{senāno sauthī bahādur sainik ghāyal thayo.} — The bravest soldier in the army was wounded.
  \item \emph{Recall:} say the Gujarati for a war, then the Gujarati for an army, then say \emph{sainik} again
\end{itemize}

\begin{tcolorbox}[breakable,colback=teal!4,colframe=teal!35!black,title={Before you move on}]
What does \gu{સૈનિક} mean? (\textquotedblleft{}A soldier\textquotedblright{}.) Say it once more.
\end{tcolorbox}
\section[\texorpdfstring{\gu{દુશ્મન} (duśman)}{દુશ્મન (duśman)}]{\gu{દુશ્મન} (duśman) — an enemy}

\label{lesson:GU-C320-dusman}



\begin{quote}
Before the new one: say the Gujarati for an army, then the Gujarati for a soldier.
\end{quote}

\subsection*{You'll want to know: \gu{દુશ્મન}}
\textbf{\gu{દુશ્મન}} — \emph{duśman} — \textquotedblleft{}an enemy\textquotedblright{}.

\begin{quote}
\gu{દુશ્મને} \gu{રાત્રે} \gu{કિલ્લો} \gu{ઘેર્યો}.
\end{quote}

\begin{quote}
\emph{duśmane rātre killo gheryo.}
\end{quote}

\begin{quote}
The enemy surrounded the fort at night.
\end{quote}

A masculine or feminine noun, by the person meant.

\begin{cousinweb}
\textbf{\gu{દુશ્મન}} (\emph{duśman}) is usually traced to Persian \emph{dušman}, \textquotedblleft{}an enemy\textquotedblright{}. As the doer it takes \textbf{\gu{એ}}: \textbf{\gu{દુશ્મને} \gu{કિલ્લો} \gu{ઘેર્યો}}.
\end{cousinweb}

\begin{grammarlens}[title={What you've built}]
One more word for telling what happened.
\end{grammarlens}

\subsection*{Your turn}
\begin{itemize}
\raggedright
  \item \emph{Say it:} \emph{duśman}
  \item \emph{Say it:} \emph{duśman}, once more
  \item \emph{Say it:} \emph{duśmane rātre killo gheryo.} — The enemy surrounded the fort at night.
  \item \emph{Recall:} say the Gujarati for an army, then the Gujarati for a soldier, then say \emph{duśman} again
\end{itemize}

\begin{tcolorbox}[breakable,colback=teal!4,colframe=teal!35!black,title={Before you move on}]
What does \gu{દુશ્મન} mean? (\textquotedblleft{}An enemy\textquotedblright{}.) Say it once more.
\end{tcolorbox}
\section[\texorpdfstring{\gu{આક્રમણ} (ākramaṇ)}{આક્રમણ (ākramaṇ)}]{\gu{આક્રમણ} (ākramaṇ) — an invasion; an attack}

\label{lesson:GU-C320-akraman}



\begin{quote}
Before the new one: say the Gujarati for a soldier, then the Gujarati for an enemy.
\end{quote}

\subsection*{You'll want to know: \gu{આક્રમણ}}
\textbf{\gu{આક્રમણ}} — \emph{ākramaṇ} — \textquotedblleft{}an invasion; an attack\textquotedblright{}.

\begin{quote}
\gu{દુશ્મને} \gu{સવારે} \gu{આક્રમણ} \gu{કર્યું}.
\end{quote}

\begin{quote}
\emph{duśmane savāre ākramaṇ karyũ.}
\end{quote}

\begin{quote}
The enemy invaded in the morning.
\end{quote}

A neuter noun.

\begin{cousinweb}
\textbf{\gu{આક્રમણ}} (\emph{ākramaṇ}) is a tatsama, Sanskrit \emph{ākramaṇa}, \textquotedblleft{}stepping upon\textquotedblright{} (CDIAL 1016), and so marching into another's land. It is done: \textbf{\gu{આક્રમણ} \gu{કર્યું}}.
\end{cousinweb}

\begin{grammarlens}[title={What you've built}]
One more word for telling what happened.
\end{grammarlens}

\subsection*{Your turn}
\begin{itemize}
\raggedright
  \item \emph{Say it:} \emph{ākramaṇ}
  \item \emph{Say it:} \emph{ākramaṇ}, once more
  \item \emph{Say it:} \emph{duśmane savāre ākramaṇ karyũ.} — The enemy invaded in the morning.
  \item \emph{Recall:} say the Gujarati for a soldier, then the Gujarati for an enemy, then say \emph{ākramaṇ} again
\end{itemize}

\begin{tcolorbox}[breakable,colback=teal!4,colframe=teal!35!black,title={Before you move on}]
What does \gu{આક્રમણ} mean? (\textquotedblleft{}An invasion; an attack\textquotedblright{}.) Say it once more.
\end{tcolorbox}
